\newcommand{\FluidBernoulliDiagram}{%
\begin{center}\begin{tikzpicture}[scale=.9,>=Latex]
 \draw[->,TextGray] (-2,-2.6)--(-2,3.4) node[above]{$z$};
 \draw[TextGray,dashed] (-2,-2.3)--(7.8,-2.3);
 \begin{scope}[rotate=16]
  \draw[thick] (0,.85)--(6,.5) (0,-.85)--(6,-.5);
  \draw[TextGray,dashed] (0,0)--(6,0);
  \draw[DeepBlue,thick] (0,0) ellipse[x radius=.20,y radius=.85];
  \draw[DeepBlue,thick] (6,0) ellipse[x radius=.117647,y radius=.5];
  \draw[TextGray,dashed] (.35,0) ellipse[x radius=.195294,y radius=.830];
  \draw[TextGray,dashed] (7.0115,0) ellipse[x radius=.117647,y radius=.5];
  \draw[->,DeepBlue,thick] (-1.1,0)--(0,0);
  \draw[->,DeepBlue,thick] (7.4,0)--(6,0);
  \node[above] at (-.9,.15) {$p_1A_1$};\node[above right] at (7.0,.6) {$p_2A_2$};
  \draw[TextGray,thin] (0,-.87)--(0,-1.28) (.35,-.87)--(.35,-1.28);
  \draw[->,TextGray] (-.55,-1.15)--(0,-1.15);
  \draw[->,TextGray] (.90,-1.15)--(.35,-1.15);
  \node[below] at (.175,-1.3) {$\dd s_1$};
  \draw[<->,TextGray] (6,-.85)--(7.0115,-.85);
  \node[below] at (6.5,-1.0) {$\dd s_2$};
  \node[above] at (0,1.02) {$A_1,\,v_1$};\node[above] at (5.9,.72) {$A_2,\,v_2$};
  \coordinate (FluidA) at (0,0);\coordinate (FluidB) at (6,0);
 \end{scope}
 \draw[<->,TextGray] (2,-2.3)--(2,0) node[midway,right]{$z_1$};
 \draw[TextGray,dotted] (FluidA)--(2,0);
 \draw[<->,TextGray] (5,-2.3)--(5,{6*sin(16)}) node[midway,left]{$z_2$};
 \draw[TextGray,dotted] (FluidB)--(5,{6*sin(16)});
\end{tikzpicture}\figcap{流管两端压力做功。相同时间内移入、移出的质量相等；高度从同一水平基准面量起。}\label{ch13:fig:bernoulli}
\end{center}}

\newcommand{\FluidShearDiagram}{%
\begin{center}\begin{tikzpicture}[scale=.85,>=Latex]
 \foreach \yy in {0,1.8}{\draw[Ink,thick] (0,\yy)--(4.2,\yy)--(5.15,\yy+.55)--(.95,\yy+.55)--cycle;}
 \foreach \yy in {.15,.60,1.05,1.50}{\draw[->,DeepBlue] (1,\yy)--({1.65+.65*\yy},\yy);}
 \draw[->,DeepBlue,very thick] (2,2.08)--(4.25,2.08);
 \node[right] at (5.5,2.08) {$v_x+\dd v_x$};
 \draw[->,DeepBlue,very thick] (2,.27)--(3.4,.27) node[right=4pt]{$v_x$};
 \draw[<->,TextGray] (-.45,0)--(-.45,1.8) node[midway,left]{$\dd y$};
 \draw[->,Ink] (-1.3,0)--(-1.3,2.7) node[above]{$y$};
 \draw[->,Ink] (0,-.8)--(2,-.8) node[right]{$x$};
 \node[above] at (1.6,2.5) {上层对下层：向右拉};
 \node[below] at (3.1,-1.05) {$\tau_{xy}=\dd F_x/\dd S=\eta\,\dd v_x/\dd y$};
\end{tikzpicture}\figcap{相距 $\dd y$ 的流层与切向速度梯度。上层流得较快，黏性作用减小两层速度差；反作用力方向相反。}\label{ch13:fig:shear}
\end{center}}

\newcommand{\FluidOutletDiagram}{%
\begin{center}\begin{tikzpicture}[scale=.78,>=Latex]
 \fill[Blue!12] (0,0) rectangle (4,3);
 \draw[thick] (0,3.6)--(0,0)--(4,0)--(4,.7) (4,1.05)--(4,3.6);
 \draw[Blue,thick] (0,3)--(4,3);
 \draw[Blue,thick] (4,.7)--(5.5,.5) (4,1.05)--(5.5,.85);
 \fill (1,3) circle(1.6pt) node[above]{$A$};
 \fill (4,.875) circle(1.6pt);
 \draw[TextGray,thin] (4,.875)--(3.4,1.3);
 \node[above left] at (3.4,1.3) {$B$};
 \draw[->,DeepBlue,thick] (4.7,.67)--(5.6,.67) node[right]{$v$};
 \draw[<->,TextGray] (-.6,.875)--(-.6,3) node[midway,left]{$h$};
 \draw[dotted,TextGray] (-.6,.875)--(3.7,.875);
 \node[above] at (2,4.0) {$p_A=p_B=p_0$};
\end{tikzpicture}\figcap{取自由液面 $A$ 与小孔出口 $B$ 比较压强、速度和高度。}\label{ch13:fig:outlet}\end{center}}

\newcommand{\FluidPlateDiagram}{%
\begin{center}\begin{tikzpicture}[scale=.85,>=Latex]
 \draw[Ink,very thick] (4,-2)--(4,2);
 \draw[Blue,thick] (0,-.25)--(3.5,-.25)..controls(3.8,-.25)and(3.8,-1)..(3.8,-1.65);
 \draw[Blue,thick] (0,.25)--(3.5,.25)..controls(3.8,.25)and(3.8,1)..(3.8,1.65);
 \draw[->,DeepBlue,thick] (.5,0)--(2,0) node[midway,above=8pt]{$v$};
 \draw[->,DeepBlue,thick] (3.65,.7)--(3.65,1.65);
 \draw[->,DeepBlue,thick] (3.65,-.7)--(3.65,-1.65);
 \draw[dashed,TextGray] (2.3,-1.75) rectangle (4.3,1.75);
 \draw[->,Ink,thick] (4.7,0)--(6,0) node[right]{$F_{\text{水对板}}$};
 \node[left] at (0,0) {$A$};\node[above] at (3.3,2.2) {虚线为控制体};
 \draw[->,TextGray] (0,-1.6)--(1.4,-1.6) node[right]{$x$};
\end{tikzpicture}\figcap{来流有正向 $x$ 动量，流出水沿板面，法向动量为零。}\label{ch13:fig:jet-plate}\end{center}}

\newcommand{\FluidSettlingDiagram}{%
\begin{center}\begin{tikzpicture}[scale=.9,>=Latex]
 \draw[thick] (0,0) circle(.62);
 \draw[->,Ink,thick] (0,-.62)--(0,-2) node[right]{$mg$};
 \draw[->,DeepBlue,thick] (-.26,{sqrt(.62*.62-.26*.26)})--(-.26,1.6) node[above left]{$F_{\!B}$};
 \draw[->,DeepBlue,thick] (.26,{sqrt(.62*.62-.26*.26)})--(.26,2.2) node[above right]{$6\pi\eta Rv_t$};
 \draw[->,TextGray,thick] (1.8,.6)--(1.8,-.7) node[right]{$v_t$};
 \node[right] at (2.2,1.6) {$F_{\!B}=\frac43\pi R^3\rho_f g$};
\end{tikzpicture}\figcap{下沉小球的受力。终端状态速度向下，浮力和黏性阻力均向上。}\label{ch13:fig:settling}\end{center}}

\section{流体运动学与物质导数}\label{ch13:sec:material}

流体由大量不断更换邻居的质元组成。拉格朗日描述（Lagrangian description）给每个质元贴上初始标签 $\vect a$，记录其位置 $\vect r(\vect a,t)$；欧拉描述（Eulerian description）则在空间各点记录速度、密度和压强场
\begin{equation}
 \vect v=\vect v(\vect r,t),\qquad \rho=\rho(\vect r,t),\qquad p=p(\vect r,t).
\end{equation}
两种描述说的是同一群质元。固定初始标签 $\vect a$，对其位置求时间导数，就得到该质元此刻的速度：
\[
 \vect v\bigl(\vect r(\vect a,t),t\bigr)
 =\left.\frac{\partial\vect r(\vect a,t)}{\partial t}\right|_{\vect a}.
\]
左边先找到质元当前所在的位置，再读取空间速度场；右边始终跟踪同一个质元。标签 $\vect a$ 不能在求导时随意改变。
竞赛流体题通常使用欧拉描述，但研究某个质元所经历的变化时必须沿它的轨迹求导。若标量场为 $f(\vect r,t)$，物质导数（material derivative）为
\begin{equation}
 \frac{\mathrm D f}{\mathrm Dt}
 =\frac{\partial f}{\partial t}+\vect v\cdot\nabla f.
\end{equation}
这是多元函数链式法则的直接结果。把沿质元轨迹读到的量写成
$f(x(t),y(t),z(t),t)$，则
\[
 \frac{\dd f}{\dd t}=\frac{\partial f}{\partial t}
 +\frac{\partial f}{\partial x}\frac{\dd x}{\dd t}
 +\frac{\partial f}{\partial y}\frac{\dd y}{\dd t}
 +\frac{\partial f}{\partial z}\frac{\dd z}{\dd t}.
\]
代入 $\dd x/\dd t=v_x$ 等关系，后三项才合写成 $\vect v\cdot\nabla f$。
第一项描述固定位置随时间的变化，后三项描述质元移动到不同位置而读到的变化；二者都不能遗漏。
因此流体质元的加速度不是单纯的 $\partial\vect v/\partial t$，而是
\begin{equation}
 \vect a=\frac{\mathrm D\vect v}{\mathrm Dt}
 =\frac{\partial\vect v}{\partial t}+(\vect v\cdot\nabla)\vect v.
\end{equation}
第二项是对流加速度：即使某点的速度场不随时间改变，质元进入速度不同的区域时仍可有加速度。

迹线是一个质元随时间走过的轨迹，满足 $\dd\vect r/\dd t=\vect v(\vect r,t)$；流线是在同一时刻处处与速度相切的曲线。定常流动中二者重合；非定常流动中一般不重合，非定常流动时不能把某一时刻画出的流线当成质元未来轨迹。
对固定时刻 $t_0$，流线可用参数 $s$ 表示为 $\dd\vect r/\dd s=\lambda(s)\vect v(\vect r,t_0)$。
在相应分量非零处，也常写成
\[
 \frac{\dd x}{v_x(x,y,z,t_0)}=\frac{\dd y}{v_y(x,y,z,t_0)}
 =\frac{\dd z}{v_z(x,y,z,t_0)}.
\]
某分量为零时应使用参数形式，不能除以零。前面的流管图中，$C_1,C_2$ 之间的侧壁由同一组流线组成，速度沿侧壁而不穿过它。

为比较相邻质元，令两位置相差 $\delta\vect r$，作一阶展开：
\[
 v_i(\vect r+\delta\vect r)=v_i(\vect r)
 +\sum_j\frac{\partial v_i}{\partial x_j}\delta x_j+O(|\delta\vect r|^2).
\]
第一项给出共同平动。剩余的速度梯度可分解成对称和反对称两部分：
\begin{equation}
 \frac{\partial v_i}{\partial x_j}
 =D_{ij}+\Omega_{ij},
 \quad
 D_{ij}=\frac12\left(\frac{\partial v_i}{\partial x_j}+\frac{\partial v_j}{\partial x_i}\right),
 \quad
 \Omega_{ij}=\frac12\left(\frac{\partial v_i}{\partial x_j}-\frac{\partial v_j}{\partial x_i}\right).
\end{equation}
$D_{ij}$ 描述局部拉伸和剪切速率，$\Omega_{ij}$ 描述刚体式转动；涡量 $\vect\omega=\nabla\times\vect v$，局部流体微团的角速度为 $\vect\omega/2$。不可压缩条件是 $\nabla\cdot\vect v=\operatorname{tr}\mathbf D=0$，表示质元体积不随运动改变，并不等价于“速度处处相同”。
具体地，$\boldsymbol\Omega\,\delta\vect r=(\vect\omega/2)\times\delta\vect r$，所以它只改变相邻质元连线的方向。
连线长度平方的变化率则为
\[
 \frac{\dd}{\dd t}|\delta\vect r|^2
 =2\,\delta\vect r\cdot\mathbf D\,\delta\vect r,
\]
反对称部分的贡献为零。故 $\mathbf D$ 才描述变形快慢；它的量纲是 $\mathrm{s^{-1}}$，不同于弹性体中无量纲的应变。

\section{静止流体}\label{ch13:sec:statics}

静止流体不能承受持续剪切应力：一旦持续受切向力，普通流体会不断变形，而不是保持静止。
任取内部一个面，其表面力只能沿法线；如果不同方向的正应力不等，换一个倾斜截面便会出现切向分量，与静止矛盾。
所以静止流体的应力各向同性。取拉伸应力为正，应力张量为
\[
 \boldsymbol\sigma=-p\mathbf I.
\]
体积 $V$ 内流体所受周围流体的压力合力为
\[
 \vect F_p=-\oint_{\partial V}p\vect n\,\dd S
 =-\int_V\nabla p\,\dd V,
\]
其中 $\vect n$ 是向外单位法向，第二步对各分量使用高斯定理。因此单位体积的压力合力为 $-\nabla p$。
若外体积力密度为 $\vect f$，静力平衡要求 $\nabla p=\vect f$。这还意味着在光滑单连通区域内
$\nabla\times\vect f=0$；注意 $\vect f$ 是单位体积的力，不能与单位质量外力混淆。
在重力场中 $\vect f=\rho\vect g$，取竖直向上为 $z$，薄层受力平衡给出
\begin{equation}
 \frac{\dd p}{\dd z}=-\rho g.
\end{equation}
不可压缩液体密度恒定时，深度增加 $h$，压强增加 $\rho gh$。同一连通静止液体中，同一水平面压强相等；
平直界面或忽略表面张力时，不同液体交界面两侧压强连续；曲面界面还可能存在后文讨论的拉普拉斯压差。
物体浸入液体后，液体对其表面的压力合力等于把相同体积液体移走所失去的压力合力：
\[
 \vect F_B=-\int_{V_{\text{排}}}\rho\vect g\,\dd V.
\]
故浮力等于排开流体的重量、方向与重力相反。密度均匀时，其作用线通过排开体积的几何中心；密度不均匀时应取排开流体的质量中心。

\subsection{从一小柱液体看压强随高度的变化}
压强不是向下的力，而是作用在所选小面上的法向力除以面积。液体内部同一点的压强是标量，受压面的朝向不同，压力方向随之改变。因此不能因为重力向下，就把侧壁受到的液体压力也画成向下。

取截面积为 $S$、底面高度为 $z$、顶面高度为 $z+\dd z$ 的竖直小柱。周围液体在底面给它向上的力 $p(z)S$，在顶面给它向下的力 $p(z+\dd z)S$；侧面压力没有竖直分量。向上为正，平衡式为
\[
 p(z)S-p(z+\dd z)S-\rho gS\dd z=0.
\]
减去相同项、除以 $S\dd z$ 并取极限，才得到 $\dd p/\dd z=-\rho g$。负号表示向上压强降低；若改用向下的深度 $h$，则 $\dd p/\dd h=\rho g$。对密度恒定的液体，从自由面 $p=p_a$ 积分到深度 $h$，有
\[
 p(h)-p_a=\int_0^h\rho g\,\dd h'=\rho gh.
\]
$p$ 是绝对压强，$p-p_a$ 是相对大气的表压。自由面与大气相通并不意味着 $p=0$，而是表压为零。比较受压面两侧的合力时，才可能把相同的大气压贡献抵消。

\subsection{不同液层与受压面的合力}
密度分层时，应在每层分别累加压强增量。在同一连通静止液体中取水平小柱，水平外体积力为零，压强差必须为零；这就是等高处等压的力学依据。若两点属于互不连通的液体，或者有水平加速度，就不能仅凭高度相同认定压强相等。

\begin{origexamplebox}{两层液体中的侧壁压力}\label{ch13:ex:teaching-layered}
一个敞口容器中，上层液体密度为 $\rho_1$、厚度为 $h_1$，下层密度为 $\rho_2>\rho_1$、厚度为 $h_2$。竖直侧壁宽度为 $b$，外侧为大气。求液体作用在侧壁下层部分上的净压力及其作用点深度，忽略界面张力。
\end{origexamplebox}
\begin{xesolution}
从两液体界面向下量 $x$，$0\le x\le h_2$。上层已贡献表压 $\rho_1gh_1$，下层再贡献 $\rho_2gx$，故
\[
 \Delta p(x)=\rho_1gh_1+\rho_2gx,
 \qquad \dd F=\Delta p(x)b\,\dd x.
\]
每一水平窄条所受净压力均沿侧壁法线向外，因而可以直接相加其大小：
\[
 F=bg\left(\rho_1h_1h_2+\frac12\rho_2h_2^2\right).
\]
求作用点不能直接取几何中心，因为下方压强更大。设合力位于界面下方 $x_c$，对界面处水平轴取矩，有
\[
 Fx_c=\int_0^{h_2}x\,\dd F
 =bg\left(\frac12\rho_1h_1h_2^2+\frac13\rho_2h_2^3\right),
\]
所以
\[
 x_c=\frac{\rho_1h_1h_2^2/2+\rho_2h_2^3/3}
 {\rho_1h_1h_2+\rho_2h_2^2/2}.
\]
距自由面的深度是 $h_1+x_c$。没有上层时 $x_c=2h_2/3$；若上层压力远大于下层压强变化，分布接近均匀，$x_c\to h_2/2$。两极限分别对应三角形和近似矩形压强分布，且 $\rho gbh^2$ 的量纲确为力。
\end{xesolution}

\subsection{浮力来自压力差，稳定性还取决于力矩}
设想把浸没物体占据的空间重新填满周围液体。静止时，这部分液体的重量被周围压力合力平衡；只要边界上的压强分布相同，把它换成物体后，压力合力也不变。这一替换论证给出浮力，而不是另外假设一种向上的力。
对于跨越两层液体的物体，应分别计算 $\rho_1gV_1+\rho_2gV_2$，不能把总排水体积都乘同一个密度。浮力与重力相等只保证竖直合力为零；物体稍微倾斜后是否产生回复力矩，还需比较重心与浮力作用线，不能仅据“浮着”判断稳定。

\begin{origexamplebox}{浮体的竖直小振动}\label{ch13:ex:teaching-buoyancy}
质量为 $m$、水平截面积恒为 $S$ 的直立柱体，在密度 $\rho$ 的液体中平衡漂浮。只允许竖直运动，忽略流体附加惯性、阻力和表面张力，且小振动中底面始终浸没、顶面始终露出。求平衡浸深及小振动周期。
\end{origexamplebox}
\begin{xesolution}
平衡时 $mg=\rho gSh_0$，所以 $h_0=m/(\rho S)$。取柱体向下的位移 $x$ 为正；此时浸深为 $h_0+x$，浮力向上且增大：
\[
 m\ddot x=mg-\rho gS(h_0+x)=-\rho gSx.
\]
平衡项消去后，剩下的正是指向平衡位置的回复力。因此
\[
 \omega^2=\frac{\rho gS}{m},\qquad
 T=2\pi\sqrt{\frac{m}{\rho gS}}.
\]
若 $x<0$，浮力减少，合力向下，也把柱体推回平衡。$\rho gS$ 的单位为 $\mathrm{N/m}$，相当于等效劲度系数。真实液体随物体一起加速会增加有效惯性，故本题结果依赖题设的忽略条件；也不能把它用于已经完全浸没、排液体积不随 $x$ 变化的柱体。
\end{xesolution}

\section{连续性方程与伯努利方程}\label{ch13:sec:bernoulli}

一般质量守恒写成
\begin{equation}
 \frac{\partial\rho}{\partial t}+\nabla\cdot(\rho\vect v)=0.
\end{equation}
对固定控制体，质量的增加率等于净流入率：
\[
 \frac{\dd}{\dd t}\int_V\rho\,\dd V
 =-\oint_{\partial V}\rho\vect v\cdot\vect n\,\dd S.
\]
把右边化成体积分，利用控制体可以任意选取，才得到上面的局部式。
再展开 $\nabla\cdot(\rho\vect v)=\vect v\cdot\nabla\rho+\rho\nabla\cdot\vect v$，可写成
\[
 \boxed{\frac{\mathrm D\rho}{\mathrm Dt}+\rho\nabla\cdot\vect v=0}.
\]
这说明散度正负控制质元体积的胀缩；质量不变时密度随之降低或升高。不可压缩流体满足
$\mathrm D\rho/\mathrm Dt=0$，因而 $\nabla\cdot\vect v=0$；不同质元原本可以有不同密度。
对稳定的一维流管，质量流率 $\dot m=\rho Av$ 沿程不变；不可压缩时化为 $Av=\text{常量}$。

无黏、不可压缩流体沿同一条流线作稳定流动时，欧拉方程（Euler equation）沿流线积分得到
\begin{equation}
 p+\frac12\rho v^2+\rho gz=\text{常量}.
\end{equation}
三项分别是单位体积的压强能、动能和重力势能。伯努利方程（Bernoulli equation）不能跨越泵、涡轮、强黏性耗散区随意使用，
也不能在旋涡流中默认所有流线具有同一常量。
下面把推导写开。对无黏质元应用牛顿定律，压力合力密度为 $-\nabla p$，故
\[
 \rho\frac{\mathrm D\vect v}{\mathrm Dt}=-\nabla p+\rho\vect g.
\]
在定常流中，沿流线取弧长 $s$，切向加速度为 $v\,\dd v/\dd s$。沿切向投影得到
\[
 \rho v\frac{\dd v}{\dd s}=-\frac{\dd p}{\dd s}-\rho g\frac{\dd z}{\dd s},
 \qquad \dd(v^2/2)+\frac{\dd p}{\rho}+g\,\dd z=0.
\]
密度为常数时逐项积分，就是伯努利方程。
\FluidBernoulliDiagram
图~\ref{ch13:fig:bernoulli}也给出功—能解释。相同短时间内，两端移过的质量相等：
$\dd m=\rho A_1\dd s_1=\rho A_2\dd s_2$。
两端压力总功为
\[
 \dd W_p=p_1A_1\dd s_1-p_2A_2\dd s_2
 =\dd m\left(\frac{p_1}{\rho}-\frac{p_2}{\rho}\right).
\]
它等于动能与重力势能的增加
$\dd m[(v_2^2-v_1^2)/2+g(z_2-z_1)]$。
出口压力做负功，正是式中压强项不能漏掉的原因。

更一般地，对定常、无黏、正压流体，若密度可写成压强的单值函数 $\rho=\rho(p)$，沿流线有
\begin{equation}
 \frac{v^2}{2}+gz+\int^p\frac{\dd p'}{\rho(p')}=\text{常量}.
\end{equation}
不可压缩时 $\rho$ 为常量，积分项化为 $p/\rho$；可压缩气流中密度和内能通常与温度耦合，不能仍把 $p+\rho v^2/2+\rho gz$ 当作普适守恒式。
若是定常、无黏、绝热且没有激波等不可逆过程的流动，单位质量内能 $u$ 满足
\[
 \dd u=-p\,\dd(1/\rho)=\frac{p}{\rho^2}\dd\rho.
\]
定义单位质量焓 $h_{\rm sp}=u+p/\rho$，则
\[
 \dd h_{\rm sp}=\dd u+\frac{\dd p}{\rho}-\frac p{\rho^2}\dd\rho
 =\frac{\dd p}{\rho}.
\]
代回沿流线的动力学式，得到
\[
 \boxed{h_{\rm sp}+\frac{v^2}{2}+gz=\text{沿该流线的常量}}.
\]
这里焓含有内能，不能只保留 $p/\rho$。例如理想气体的内能随温度改变，气流加速可能伴随温度降低。
受热驱动的自然对流还涉及密度变化和浮力；液体密度变化虽小，也不能一概删去其浮力效应。
常用近似是在质量守恒中取无散速度场，而在重力项中保留小密度差；这与“所有动力学都与温度无关”不同。

\begin{origexamplebox}{小孔流出速度}\label{ch13:ex:outlet}
大水箱自由液面到侧壁小孔的高度差为 $h$，水箱横截面远大于小孔。忽略黏性，求流出速度。
\end{origexamplebox}
\begin{xesolution}
在自由液面与小孔出口之间使用伯努利方程。两点都与大气相通，压强项相消；大水箱自由液面速度近似为零：
\FluidOutletDiagram
图~\ref{ch13:fig:outlet}中 $z_A-z_B=h$，$v_A\simeq0$，于是
\[
 p_0+\rho gh=p_0+\frac12\rho v^2.
\]
所以
\[
 \boxed{v=\sqrt{2gh}}.
\]
这与物体自由下落高度 $h$ 的速度相同，但流体微元的真实轨迹和压强变化并不等同于一个孤立质点自由落体。
\end{xesolution}

\section{动量流与流体作用力}\label{ch13:sec:momentum}

稳定流动中，控制体所受外力等于流出动量流率减去流入动量流率：
\begin{equation}
 \sum\vect F=\sum_{\rm out}\dot m\vect v-
 \sum_{\rm in}\dot m\vect v.
\end{equation}
这里的外力包括压强、重力和固体壁面对流体的力。求流体对管壁的力时，最后还要用牛顿第三定律反向。

\begin{origexamplebox}{水射流正撞平板}\label{ch13:ex:jet-plate}
密度为 $\rho$、截面积为 $A$、速度为 $v$ 的稳定水射流垂直撞击静止大平板，撞后水沿板面向四周流走。忽略重力，求水对板的平均力。
\end{origexamplebox}
\begin{xesolution}
取法向为 $x$ 轴。质量流率 $\dot m=\rho Av$。流入法向速度为 $v$，流出水流沿板面，法向速度为零。
\FluidPlateDiagram
图~\ref{ch13:fig:jet-plate}中虚线包围转向区；图右箭头是水对板的力，列控制体方程时应使用它的反作用力。
平板对水的力为
\[
 F_{\text{板对水},x}=\dot m(0-v)=-\rho Av^2.
\]
所以水对板的力方向沿来流方向，大小
\[
 \boxed{F=\rho Av^2}.
\]
量纲为 $\mathrm{kg\,m^{-3}}\cdot\mathrm{m^2}\cdot\mathrm{m^2s^{-2}}=\mathrm N$。
\end{xesolution}

\subsection{为什么稳态流动仍然需要外力}
稳态表示固定位置的流速不随时间变化，不表示经过这里的每一份流体都没有加速度。图~\ref{ch13:fig:jet-plate} 中，虚线区域内的水总量可以不变，但新进入的水不断丢失法向动量，这就需要平板持续施力。
在短时间 $\dd t$ 内，流入质量为 $\dot m\dd t$，流入动量为 $\dot m\vect v_{\rm in}\dd t$；流出端同理。稳态时控制体内的动量不积累，所以外力冲量必须补足出入动量之差。除以 $\dd t$ 得到上面的动量流方程。多出口时必须先按矢量相加，不能直接相加速率。

对管内流动，在入口截面，外部流体把所研究流体推入控制体；在出口截面，外部流体阻碍它流出。两个压力方向均指向控制体内部。若求管道与环境之间的净作用力，使用表压时还须把管外大气压的抵消计入，不能只把绝对压强随意置零。

\begin{origexamplebox}{射流推动匀速后退的平板}\label{ch13:ex:teaching-moving-plate}
沿水平 $+x$ 方向的细水射流速度为 $v$、截面积为 $A$、密度为 $\rho$。一块足够大的平板沿同方向以恒速 $u$ 后退，$0\le u<v$。水撞板后沿板面流走，相对平板没有法向速度。忽略重力及空气作用，求水对板的力和传给板的功率。
\end{origexamplebox}
\begin{xesolution}
以匀速平板为参考系，它仍是惯性系。对应图~\ref{ch13:fig:jet-plate} 的控制体随板平移，入射相对速度为 $v-u$，故每秒真正撞板的质量为
\[
 \dot m_{\rm hit}=\rho A(v-u).
\]
这不是固定截面的 $\rho Av$：一部分水还在追赶平板。水相对板的法向速度从 $v-u$ 变成零，每单位质量丢失法向动量 $v-u$。由反作用关系，板受到沿 $+x$ 的力
\[
 F=\dot m_{\rm hit}(v-u)=\rho A(v-u)^2.
\]
也可在地面系看同一份水，撞前法向速度为 $v$，撞后为 $u$，其动量变化给出同一个结果。力对平板的功率必须使用地面系板速：
\[
 P=Fu=\rho Au(v-u)^2.
\]
在 $0<u<v$ 内，$\dd P/\dd u=\rho A(v-u)(v-3u)$，故最大功率出现在 $u=v/3$，大小为 $4\rho Av^3/27$。$u=0$ 时有力却不向静止板传机械功；$u\to v$ 时几乎没有水追上平板，力和功率均趋零。这也说明力最大与功率最大不是同一个条件。
\end{xesolution}

\section{黏性、层流与相似准则}\label{ch13:sec:viscosity}

牛顿流体（Newtonian fluid）的剪切应力与速度梯度成正比：
\begin{equation}
 \tau=\eta\frac{\partial v_x}{\partial y},
\end{equation}
其中 $\eta$ 为动力黏度。半径 $R$、长度 $L$ 的圆管内充分发展层流满足哈根--泊肃叶定律（Hagen--Poiseuille law）
\begin{equation}
 Q=\frac{\pi R^4}{8\eta L}\Delta p.
\end{equation}
\FluidShearDiagram
图~\ref{ch13:fig:shear}中，上下两流层沿 $x$ 流动，速度只随法向坐标 $y$ 改变。
速度差除以层间距是剪切速率；乘以动力黏度才是单位面积的切向力。
$\eta$ 的单位为 $\mathrm{Pa\,s}$，另定义运动黏度
\[
 \boxed{\nu=\eta/\rho},\qquad [\nu]=\mathrm{m^2/s}.
\]
推导圆管流量时，设 $G=\Delta p/L>0$，距轴 $r$ 处的定常速度为 $v_z(r)$。沿轴的力平衡给出
$\eta r^{-1}\dd[r\,\dd v_z/\dd r]/\dd r=-G$。
积分一次为 $r\,\dd v_z/\dd r=-Gr^2/(2\eta)+C$；轴心速度梯度有限要求 $C=0$。
再积分并用管壁无滑移条件 $v_z(R)=0$：
\[
 v_z(r)=\frac{G}{4\eta}(R^2-r^2),\qquad
 Q=\int_0^R v_z(r)\,2\pi r\dd r
 =\frac{\pi GR^4}{8\eta}.
\]
这一步同时说明充分发展、直圆管、不可压缩牛顿流体和无滑移边界等条件为何重要。
流量对半径的四次方依赖说明细管尺寸误差会产生很大影响。雷诺数（Reynolds number）
\begin{equation}
 \mathrm{Re}=\frac{\rho v\ell}{\eta}
\end{equation}
比较惯性与黏性效应；其临界值依赖几何和扰动，不能把某个数值当成所有流动的普适分界。

低雷诺数下，半径 $R$ 的球以速度 $v$ 在黏性流体中缓慢运动，斯托克斯阻力（Stokes drag）为 $6\pi\eta Rv$。

\begin{origexamplebox}{小球在黏性液体中的终端速度}\label{ch13:ex:settling}
密度为 $\rho_s$、半径为 $R$ 的小球在密度 $\rho_f$、黏度 $\eta$ 的液体中下落。假设满足斯托克斯阻力，求终端速度。
\end{origexamplebox}
\begin{xesolution}
终端状态合力为零。向下重力减去浮力等于向上黏性阻力：
\FluidSettlingDiagram
图~\ref{ch13:fig:settling}按下沉情形画出三个力；向下为正，平衡式为
\[
 \frac43\pi R^3(\rho_s-\rho_f)g=6\pi\eta Rv_t.
\]
因此
\[
 \boxed{v_t=\frac{2(\rho_s-\rho_f)gR^2}{9\eta}}.
\]
该式要求 $\rho_s>\rho_f$ 且雷诺数足够小；若小球上浮，速度方向反向，大小用密度差绝对值。
解出后可以代回 $\mathrm{Re}=2\rho_fR|v_t|/\eta$ 检查是否远小于1。
这里特征长度取球直径 $2R$；若检验不满足，就必须更换阻力模型，而不是继续相信代数上求出的速度。
\end{xesolution}

\section{表面张力}\label{ch13:sec:surface}

表面张力 $\gamma$ 既可理解为沿界面单位长度的拉力，也可理解为增加单位面积界面所需的可逆功。
曲面两侧压强差满足Laplace公式；球形液滴只有一个界面，肥皂泡有内外两个界面：
\begin{equation}
 \Delta p_{\rm drop}=\frac{2\gamma}{R},\qquad
 \Delta p_{\rm bubble}=\frac{4\gamma}{R}.
\end{equation}
半径越小，维持曲率所需压强差越大。毛细管中接触角为 $\theta$ 时，液面升降高度为
\begin{equation}
 h=\frac{2\gamma\cos\theta}{\rho gr}.
\end{equation}

\subsection{同一个系数为什么既是力除以长度，又是功除以面积}
设薄液膜有一条可滑动的直边，长度为 $\ell$，准静态向外移动 $\dd x$。液膜有前后两个界面，面积总增量为 $\dd A=2\ell\dd x$，外力做功为 $F\dd x$。在温度、组成不变，且过程可逆的条件下，$F\dd x=\gamma\dd A$，所以 $F=2\gamma\ell$。若只有一个界面，相应的因子2消失。
因此 $\mathrm{N/m}=\mathrm{J/m^2}$ 并非偶然的单位巧合，而是力做功与面积变化之间的关系。液膜沿界面拉紧边界，不是处处沿固定空间方向施力。

\subsection{球面压差与毛细升降的推导}
先看单界面的球形液滴。令半径准静态增加 $\dd R$，体积和面积的变化为
\[
 \dd V=4\pi R^2\dd R,\qquad \dd A=8\pi R\dd R.
\]
内外压强差做功为 $(p_{\rm in}-p_{\rm out})\dd V$，用于增加界面自由能 $\gamma\dd A$。约去公共因子后，$p_{\rm in}-p_{\rm out}=2\gamma/R$。薄壁肥皂泡有两个近似同半径的界面，面积变化加倍，压差才成为 $4\gamma/R$；“两个界面”不是把半径也乘2。

再看半径为 $r$ 的竖直毛细管，管外液面宽广。接触角 $\theta$ 从液体内部量取，圆形接触线长 $2\pi r$，界面张力的竖直合力为 $2\pi r\gamma\cos\theta$。在弯月面体积相对液柱体积可忽略、气体密度可忽略的近似下，它平衡高出外液面的液柱重量：
\[
 2\pi r\gamma\cos\theta=\rho g\pi r^2h.
\]
于是得到前式。$h$ 是相对管外液面的有符号高度：润湿时 $\theta<\pi/2$、$h>0$；不润湿时 $\theta>\pi/2$、$h<0$，表示液面下降。忽略弯月面体积通常要求管半径远小于毛细长度 $\sqrt{\gamma/(\rho g)}$。接触角变化、管壁污染或液柱尚在运动时，不能原样套用静态结果。

\begin{origexamplebox}{液滴合并时的界面能变化}\label{ch13:ex:teaching-droplets}
两滴同种不可压缩液体均为半径 $R$ 的球形液滴，合并后形成一个球形液滴。设表面张力为常量 $\gamma$，忽略重力。求合并后的半径、界面能减少量及新液滴的内外压差。界面能的减少是否等于总能量的消失？
\end{origexamplebox}
\begin{xesolution}
先守恒的是体积而不是表面积。设新半径为 $R_f$，则
\[
 \frac43\pi R_f^3=2\frac43\pi R^3,
 \qquad R_f=2^{1/3}R.
\]
初、末总表面积分别为 $8\pi R^2$ 和 $4\pi 2^{2/3}R^2$。因此界面能减少
\[
 \Delta E_s=4\pi\gamma R^2(2-2^{2/3})>0.
\]
这部分能量可先转化为液滴内部流动、形变振动的能量，再由黏性耗散转化为内能；没有凭空消失。若过程中温度明显改变，常量 $\gamma$ 的近似还须重查。
新液滴压差为 $\Delta p_f=2\gamma/(2^{1/3}R)$，比原压差小。半径增大、曲率减小与压差下降一致；$\gamma R^2$ 的量纲是能量，$\gamma/R$ 的量纲是压强。
\end{xesolution}

\section{流体运动方程（选读）}\label{ch13:sec:equations-reading}
\begin{starread}{欧拉方程与纳维--斯托克斯方程}
无黏流体的局部动力学方程为
\[
 \rho\left(\frac{\partial\vect v}{\partial t}+\vect v\cdot\nabla\vect v\right)
 =-\nabla p+\rho\vect g.
\]
加入牛顿黏性应力后得到不可压缩纳维--斯托克斯方程（Navier--Stokes equation）
\[
 \rho\left(\frac{\partial\vect v}{\partial t}+\vect v\cdot\nabla\vect v\right)
 =-\nabla p+\eta\nabla^2\vect v+\rho\vect g,
 \qquad \nabla\cdot\vect v=0.
\]
非线性对流项使涡旋、边界层和湍流问题远比一维管流复杂。竞赛计算若采用伯努利或泊肃叶模型，
必须先说明忽略或保留了哪些项。

对可压缩牛顿流体，黏性应力的常用形式为
\[
 \boldsymbol\tau=\eta\left[\nabla\vect v+(\nabla\vect v)^{\mathsf T}
 -\frac23(\nabla\cdot\vect v)\mathbf 1\right]+\zeta(\nabla\cdot\vect v)\mathbf 1,
\]
其中 $\eta$ 是剪切黏度，$\zeta$ 是体黏度。不可压缩且 $\eta$ 均匀时才简化为动量方程中的 $\eta\nabla^2\vect v$。这说明泊肃叶流、低雷诺数绕流和可压缩声流采用的黏性模型虽同源，保留的项并不相同。
其来由是把总应力分为 $\boldsymbol\sigma=-p\mathbf I+\boldsymbol\tau$。
不含局部力偶的微小流体立方体满足角动量平衡，要求 $\tau_{xy}=\tau_{yx}$ 等，因此应力张量对称。
形变率还可分成无迹部分与纯体积变化：
\[
 \mathbf D=\left(\mathbf D-\frac13\nabla\cdot\vect v\,\mathbf I\right)
 +\frac13\nabla\cdot\vect v\,\mathbf I.
\]
各向同性牛顿流体的应力分别与这两部分成正比，所以出现剪切黏度和体黏度两个系数。
热力学压强 $p$ 与总正应力的平均值在可压缩黏性流动中未必相同，不能把包含体黏度项的全部 $\boldsymbol\tau$ 都称为无迹偏应力。
对任意小控制体，表面应力的合力密度为 $\nabla\cdot\boldsymbol\sigma$，故一般动量式为
\[
 \rho\frac{\mathrm D\vect v}{\mathrm Dt}=-\nabla p+\nabla\cdot\boldsymbol\tau+\rho\vect g.
\]
若 $\eta,\zeta$ 均匀，逐项求散度得到
\[
 \nabla\cdot\boldsymbol\tau=\eta\nabla^2\vect v+
 \left(\zeta+\frac\eta3\right)\nabla(\nabla\cdot\vect v).
\]
再用不可压缩条件，后一项才消失。方程虽然形式紧凑，三维非线性流动的存在性、正则性以及湍流描述仍涉及深刻数学问题；写出方程不等于已经求出了任意流动。
\end{starread}
